\documentclass{beamer}
\usepackage{setspace}
\usepackage{bm,color,float,epstopdf,amsthm,epstopdf}
\usepackage{multirow,tikz}
\newcommand{\fttheory}[1]{\fcolorbox{blue}{blue}{\color{white}\textbf{#1}\hspace{\textwidth}}}
\newcommand{\ftexampl}[1]{\fcolorbox{orange}{orange}{\color{black}\textbf{#1}\hspace{\textwidth}}}
\newcommand{\ftdata}[1]{\fcolorbox{violet}{violet}{\color{white}\textbf{#1}\hspace{\textwidth}}}
\newcommand{\ftheader}[1]{\fcolorbox{lightgray}{lightgray}{\color{black}\textbf{#1}\hspace{\textwidth}}}
\newcommand{\ftvocab}[1]{\fcolorbox{cyan}{cyan}{\color{black}\textbf{#1}\hspace{\textwidth}}}
\newcommand{\usd}{\text{ {\scriptsize USD}}}
\newcommand{\eur}{\text{ {\scriptsize EUR}}}
\newcommand{\aud}{\text{ {\scriptsize AUD}}}
\newcommand{\mxn}{\text{ {\scriptsize MXN}}}
\newcommand{\musd}{\text{ m {\scriptsize USD}}}
\newcommand{\meur}{\text{ m {\scriptsize EUR}}}
\newcommand{\maud}{\text{ m {\scriptsize AUD}}}

\renewcommand{\arraystretch}{1.5}
\begin{document}

% Factors Affecting Option Prices
\begin{frame}{\fttheory{Factors Affecting Option Prices}}
	\begin{itemize}
		\item \textbf{Strike Price and Stock Price:} the value of a call is \underline{decreasing} in the strike price (the higher the strike, the lower the value); the value of a call is \underline{increasing} in the price of the stock (the higher the price of the stock, the higher the value). 
		\item \textbf{Arbitrage Bounds on Option Prices:} 
		\begin{itemize}
			\item an American option carries all of the same rights and privileges as an otherwise equivalent European option. The American option can't be worth less. 
			\item a put option can't be worth more than its strike price and a call can't be worth more than the stock itself.
		\end{itemize}
	\end{itemize}
\end{frame}

% Factors Affecting Option Prices
\begin{frame}{\fttheory{Factors Affecting Option Prices}}
	\begin{itemize}
		\item \textbf{Option Prices and the Exercise Date:} 
		\begin{itemize}
			\item (for American options,) the longer the time to the exercise date, the more valuable the option.
			\item An American option with a later exercise date cannot be worth less than an otherwise identical American option with an earlier exercise date.
			\item These arguments \underline{do not} carry through for European options.
		\end{itemize}
		\item \textbf{Option Prices and Volatility:} 
		\begin{itemize}
			\item The value of an option generally increases with the volatility of the stock.
			\item Increasing the volatility of the stock increases the likelihood that the stock will take on extreme values (which is where options are valuable).
		\end{itemize}
	\end{itemize}
\end{frame}

\end{document}
